\noindent\textbf{Elección del procedimiento.} Antes de aplicar Hagen--Poiseuille hay que calcular la velocidad media y comprobar el régimen con Reynolds. Para flujo laminar plenamente desarrollado, $\Delta p=128\mu LQ/(\pi D^4)$; la potencia hidráulica perdida es $\Delta p\,Q$.

\begin{enumerate}

\item El dato es diámetro, por lo que el área circular lleva un factor $1/4$.

\[\fitmath{A=\frac{\pi\cdot {D}^{2}}{4}}\]

\[\fitmath{A=\frac{\pi\cdot {(0.006)}^{2}}{4}\approx 2.827433\times10^{-5}\,\mathrm{m^{2}}}\]

\item De continuidad $Q=AV$, despejar la velocidad media.

\[\fitmath{V=\frac{Q}{A}}\]

\[\fitmath{V=\frac{(0.00015)}{(2.827433388\times10^{-5})}\approx 5.305165\,\mathrm{m/s}}\]

\item Comparar inercia y viscosidad con la escala diámetro.

\[\fitmath{\mathrm{Re}=\frac{\rho\cdot V\cdot D}{\mu}}\]

\[\fitmath{\mathrm{Re}=\frac{(850)\cdot (5.30516477)\cdot (0.006)}{(0.02)}\approx 1352.817\,}\]

\item Con la estimación de entrada $L_e/D\simeq0.0575\mathrm{Re}$ usada en este capítulo, comprobar si $L_e$ es pequeño frente a $L$.

\[\fitmath{\mathrm{Le}=0.0575\cdot \mathrm{Re}\cdot D}\]

\[\fitmath{\mathrm{Le}=0.0575\cdot (1352.817016)\cdot (0.006)\approx 0.4667219\,\mathrm{m}}\]

\item Despejar la pérdida de presión de la relación caudal--presión laminar.

\[\fitmath{\mathrm{dp}=\frac{128\cdot \mu\cdot L\cdot Q}{\pi\cdot {D}^{4}}}\]

\[\fitmath{\mathrm{dp}=\frac{128\cdot (0.02)\cdot (3)\cdot (0.00015)}{\pi\cdot {(0.006)}^{4}}\approx 282942.1\,\mathrm{Pa}}\]

\item Multiplicar caída de presión por caudal, no por velocidad.

\[\fitmath{P=\mathrm{dp}\cdot Q}\]

\[\fitmath{P=(282942.1211)\cdot (0.00015)\approx 42.44132\,\mathrm{W}}\]

\end{enumerate}

\noindent\textbf{Contraste y comparación.} La estimación de entrada da aproximadamente $0.467$ m frente a 3 m. La relación de tramo plenamente desarrollado omite la pérdida adicional de entrada; su uso aquí es una aproximación, no una solución exacta de todo el conducto.
